\chapter{Linear Combinations, Span, and Geometry}

In the last chapter you met a vector as three things at once: a list of numbers, a point in space, and an arrow from the origin. That triple identity is the whole game. Now we get to the single most important verb in all of linear algebra. Not addition. Not multiplication. \emph{Combination.}

Almost everything in this subject is some version of one question: \emph{starting from a few vectors, where can I get to by scaling them and adding them up?} The set of reachable destinations has a name -- the \textbf{span} -- and the shapes it takes (a point, a line, a plane, all of space) are the geometry hiding inside every system of equations, every matrix, every dataset. If you really feel what a span looks like, half of Strang's course will feel like remembering rather than learning.

So slow down here. This chapter is short on machinery and long on pictures. Let us build it.

\section{Scaling and Adding: The Two Moves}

You only ever do two things to vectors in linear algebra, and you met both already. Let us name them out loud because the entire chapter is built from exactly these two moves and nothing else.

\textbf{Move one: scaling.} Take a vector $\vec{v}$ and a number $c$ (a \emph{scalar}). Then $c\,\vec{v}$ stretches $\vec{v}$ by the factor $c$. If $c>1$ it grows; if $0<c<1$ it shrinks; if $c<0$ it flips to point the other way; if $c=0$ it collapses to the origin. Crucially, scaling never lets $\vec{v}$ leave its own line. No matter what $c$ you pick, $c\,\vec{v}$ sits somewhere on the infinite line through the origin and the tip of $\vec{v}$.

\textbf{Move two: adding.} Take two vectors and add them tip-to-tail (the parallelogram rule from the last chapter). Addition lets you leave a single line and explore sideways.

\begin{intuitionbox}
Scaling moves you \emph{along} a direction. Adding lets you \emph{mix} directions. A linear combination is just: pick how much of each direction you want (scaling), then mix them (adding). That is it. The rest of this chapter is watching what shapes you can reach by doing this.
\end{intuitionbox}

\begin{figure}[H]
\centering
\begin{tikzpicture}[scale=1.0,>=Stealth]
  \draw[help lines,gray!30] (-3.4,-1.4) grid (3.4,3.4);
  \draw[->] (-3.4,0)--(3.4,0) node[right]{$x$};
  \draw[->] (0,-1.4)--(0,3.4) node[above]{$y$};
  % base vector v = (1,1)
  \draw[->,vecblue,very thick] (0,0)--(1,1) node[midway,above left]{$\vec{v}$};
  % scaled copies
  \draw[->,vecblue,thick,dashed] (0,0)--(2,2) node[right]{$2\vec{v}$};
  \draw[->,vecblue,thick,dashed] (0,0)--(-1,-1) node[below left]{$-\vec{v}$};
  % the whole line
  \draw[vecblue!50,thick] (-1.3,-1.3)--(3.3,3.3);
  \node[vecblue!70!black] at (2.6,1.7) {\small the line $c\vec{v}$};
\end{tikzpicture}
\caption{Scaling alone. Every scalar multiple $c\,\vec{v}$ of a single nonzero vector lands on one line through the origin. Scaling can stretch, shrink, or flip, but it cannot escape the line.}
\label{ch03:scaling-line}
\end{figure}

\section{The Linear Combination}

Now we combine the two moves into one definition. Give me vectors $\vec{v}_1, \vec{v}_2, \dots, \vec{v}_k$ and scalars $c_1, c_2, \dots, c_k$. The expression
\[
  c_1 \vec{v}_1 + c_2 \vec{v}_2 + \cdots + c_k \vec{v}_k
\]
is called a \textbf{linear combination} of the vectors. The scalars $c_i$ are the \textbf{weights} (or coefficients). The word ``linear'' is doing real work: the only operations allowed are scaling (linear in each weight) and adding. No squaring a vector, no multiplying two vectors together, no taking a sine. Just scale and add.

\begin{definition}[Linear combination]\label{ch03:def-lincomb}
A vector $\vec{b}$ is a \textbf{linear combination} of $\vec{v}_1,\dots,\vec{v}_k \in \R^n$ if there exist scalars $c_1,\dots,c_k \in \R$ such that
\[
  \vec{b} = c_1\vec{v}_1 + c_2\vec{v}_2 + \cdots + c_k\vec{v}_k.
\]
\end{definition}

Let us make it utterly concrete. Take two vectors in the plane,
\[
  \vec{v}_1 = \colvec{2\\1}, \qquad \vec{v}_2 = \colvec{1\\3}.
\]
Here is the linear combination with weights $c_1 = 2$ and $c_2 = 1$:
\[
  2\colvec{2\\1} + 1\colvec{1\\3}
  = \colvec{4\\2} + \colvec{1\\3}
  = \colvec{5\\5}.
\]
And with weights $c_1 = 1$, $c_2 = -1$:
\[
  1\colvec{2\\1} - 1\colvec{1\\3}
  = \colvec{2\\1} - \colvec{1\\3}
  = \colvec{1\\-2}.
\]
Each choice of weights gives a destination. Different weights, different destination. The grand question of the chapter is: as we sweep the weights over \emph{all} possible real numbers, what is the full set of destinations?

\begin{figure}[H]
\centering
\begin{tikzpicture}[scale=0.75,>=Stealth]
  \draw[help lines,gray!25] (-1.4,-2.4) grid (5.9,5.9);
  \draw[->] (-1.4,0)--(5.9,0) node[right]{$x$};
  \draw[->] (0,-2.4)--(0,5.9) node[above]{$y$};
  % v1, v2
  \draw[->,vecblue,very thick] (0,0)--(2,1) node[midway,below right]{$\vec{v}_1$};
  \draw[->,vecred,very thick] (0,0)--(1,3) node[midway,above left]{$\vec{v}_2$};
  % build 2v1 + v2
  \draw[->,vecblue,thick,dashed] (0,0)--(4,2) node[midway,below right]{$2\vec{v}_1$};
  \draw[->,vecred,thick,dashed] (4,2)--(5,5);
  \draw[->,targetgreen,very thick] (0,0)--(5,5) node[above right]{$2\vec{v}_1+\vec{v}_2$};
  \fill[targetgreen] (5,5) circle (2.5pt);
\end{tikzpicture}
\caption{The linear combination $2\vec{v}_1 + \vec{v}_2 = \colvec{5\\5}$ built tip-to-tail: walk out $2\vec{v}_1$ (blue, doubled), then add $\vec{v}_2$ (red) from there. The green arrow is the result.}
\label{ch03:one-combo}
\end{figure}

\begin{focusbox}
A linear combination is a recipe: ``so much of this direction, so much of that direction, add them up.'' Hold onto the recipe metaphor. The weights are the amounts of each ingredient; the result is the dish. The span (coming next) is the menu of every dish you can possibly cook.
\end{focusbox}

\section{Span: The Set of All Combinations}

Here is the star of the chapter.

\begin{definition}[Span]\label{ch03:def-span}
The \textbf{span} of vectors $\vec{v}_1,\dots,\vec{v}_k$, written $\spn\{\vec{v}_1,\dots,\vec{v}_k\}$, is the set of \emph{all} their linear combinations:
\[
  \spn\{\vec{v}_1,\dots,\vec{v}_k\}
  = \set{\, c_1\vec{v}_1 + \cdots + c_k\vec{v}_k : c_1,\dots,c_k \in \R \,}.
\]
\end{definition}

Read that as a verb and a noun. As a verb: the vectors \emph{span} a region -- they reach it, they cover it. As a noun: the span \emph{is} that region, the complete set of reachable points. The weights are the dials; the span is everywhere the dials can take you.

The beautiful fact is that the span is never a random scatter of points. It is always one of a very short list of clean geometric shapes: the origin alone, a line, a plane, all of $3$-space, and so on up. Let us walk up that ladder one rung at a time, because the ladder \emph{is} the geometry of linear algebra.

\subsection{Span of one vector: a line (usually)}

Start with a single nonzero vector $\vec{v}$. Its span is the set of all $c\,\vec{v}$. We already saw this in Figure~\ref{ch03:scaling-line}: scaling one vector traces out a line through the origin. So
\[
  \spn\{\vec{v}\} = \text{the line through the origin in the direction of } \vec{v}.
\]

There is one exception, and it matters. If $\vec{v} = \vec{0}$, then $c\,\vec{0} = \vec{0}$ for every $c$. The span of the zero vector is just the single point $\{\vec{0}\}$ -- a very lonely span. Scaling zero gets you nowhere.

\begin{trapbox}
``A line'' here always means a line \emph{through the origin}. The span of one vector can never be a line that misses the origin, because $c=0$ always produces $\vec{0}$, so the origin is in \emph{every} span. If someone hands you a line that does not pass through the origin and asks ``is this a span?'' the answer is no. (It is a \emph{shifted} line -- an affine line -- and we will meet those when we solve $\mat{A}\vec{x}=\vec{b}$ for nonzero $\vec{b}$.)
\end{trapbox}

\subsection{Span of two vectors: a plane (usually)}

Now take two vectors $\vec{v}_1$ and $\vec{v}_2$ and ask for $\spn\{\vec{v}_1,\vec{v}_2\} = \set{c_1\vec{v}_1 + c_2\vec{v}_2}$.

Think of it as a grid. Sliding $c_1$ moves you along the $\vec{v}_1$ direction; sliding $c_2$ moves you along the $\vec{v}_2$ direction. With two independent dials you can reach any point on a flat sheet -- a \textbf{plane} through the origin. The two vectors lay down a (possibly skewed) coordinate grid, and the span is the whole infinite sheet that grid lives on.

\begin{figure}[H]
\centering
\begin{tikzpicture}[scale=0.85,>=Stealth]
  \draw[help lines,gray!20] (-2.4,-2.4) grid (4.4,4.4);
  \draw[->] (-2.4,0)--(4.4,0) node[right]{$x$};
  \draw[->] (0,-2.4)--(0,4.4) node[above]{$y$};
  % skewed grid from v1=(2,1), v2=(1,2)
  \foreach \c in {-2,-1,0,1,2}{
    \draw[vecblue!30] ($\c*(2,1)+(-2,-4)$)--($\c*(2,1)+(2,4)$);
    \draw[vecred!30]  ($\c*(1,2)+(-4,-2)$)--($\c*(1,2)+(4,2)$);
  }
  \draw[->,vecblue,very thick] (0,0)--(2,1) node[below right]{$\vec{v}_1$};
  \draw[->,vecred,very thick]  (0,0)--(1,2) node[above left]{$\vec{v}_2$};
  \fill[targetgreen] (3,3) circle (2.5pt) node[above right,targetgreen]{\small $\vec{v}_1+\vec{v}_2$};
\end{tikzpicture}
\caption{Two non-parallel vectors lay down a skewed grid. Sliding the two weights $c_1,c_2$ reaches every point of the plane -- here, in $\R^2$, that plane is the entire flat page. The blue and red families of lines are the ``coordinate lines'' of the grid.}
\label{ch03:two-grid}
\end{figure}

In $\R^2$ that plane \emph{is} all of $\R^2$: two non-parallel vectors in the plane span the whole plane. (We will prove this cleanly below.) But the same two vectors living in $\R^3$ span a genuine $2$-dimensional plane floating inside $3$-space -- a flat sheet through the origin, with all of $3$-space stretching above and below it that you cannot reach.

\textbf{The exception again.} What if $\vec{v}_2$ happens to lie on the same line as $\vec{v}_1$ -- say $\vec{v}_2 = 3\vec{v}_1$? Then the second dial gives you nothing new. Anything you build, $c_1\vec{v}_1 + c_2(3\vec{v}_1) = (c_1 + 3c_2)\vec{v}_1$, is still just a multiple of $\vec{v}_1$. The span collapses back down to a single line. Two vectors span a plane only when they point in genuinely different directions.

\subsection{Span of three vectors: all of space (usually)}

Three vectors in $\R^3$ that point in three genuinely different directions span all of $\R^3$. Three dials, three independent directions, and you can reach every point in space.

But watch the exceptions stack up, because they are the whole story of \emph{dependence}:

\begin{itemize}
  \item If the third vector $\vec{v}_3$ already lies in the plane spanned by $\vec{v}_1$ and $\vec{v}_2$, it adds no new direction. The span stays a plane.
  \item If all three lie on one line, the span is just that line.
  \item If all three are $\vec{0}$, the span is the origin.
\end{itemize}

\begin{focusbox}
The ladder of spans: \emph{point, line, plane, space, \dots} Each new vector either bumps you up one rung (it points somewhere genuinely new) or it does nothing (it was already reachable). That binary -- ``new direction'' versus ``nothing new'' -- is the seed of \emph{linear independence}, the topic of a coming chapter. You are meeting it here, geometrically, before its name.
\end{focusbox}

\begin{figure}[H]
\centering
\begin{tikzpicture}[scale=1.5,>=Stealth,
   x={(0.9cm,0.32cm)}, y={(-0.86cm,0.42cm)}, z={(0cm,1cm)}]
  % a plane through origin spanned by v1, v2, with a third vector poking out
  \fill[targetgreen!12] (-1.2,-1.2,0)--(1.4,-1.2,0)--(1.4,1.2,0)--(-1.2,1.2,0)--cycle;
  \draw[gray!50] (-1.2,-1.2,0)--(1.4,-1.2,0)--(1.4,1.2,0)--(-1.2,1.2,0)--cycle;
  % axes hint
  \draw[->,gray] (0,0,0)--(1.6,0,0) node[right]{\small $x$};
  \draw[->,gray] (0,0,0)--(0,1.5,0) node[left]{\small $y$};
  \draw[->,gray] (0,0,0)--(0,0,1.5) node[above]{\small $z$};
  % v1, v2 in plane
  \draw[->,vecblue,very thick] (0,0,0)--(1,0.3,0) node[right]{$\vec{v}_1$};
  \draw[->,vecred,very thick]  (0,0,0)--(0.2,1,0) node[left]{$\vec{v}_2$};
  % v3 out of plane
  \draw[->,structurepurple,very thick] (0,0,0)--(0.3,0.2,1) node[above]{$\vec{v}_3$};
  \node[targetgreen!60!black] at (-0.7,0.9,0) {\small $\spn\{\vec{v}_1,\vec{v}_2\}$};
\end{tikzpicture}
\caption{In $\R^3$: $\vec{v}_1,\vec{v}_2$ span the shaded plane through the origin. A third vector $\vec{v}_3$ that pokes \emph{out} of that plane is a brand new direction -- together the three span all of $\R^3$. If $\vec{v}_3$ had lain flat in the plane, it would have added nothing.}
\label{ch03:three-3d}
\end{figure}

\begin{walkbox}
Picture a single vector in $\R^3$, drawn as an arrow. Its span is a line. Now picture sweeping that line around -- how many genuinely \emph{different} lines through the origin are there? Infinitely many, one for each direction. Now: how many genuinely different \emph{planes} through the origin are there in $\R^3$? Also infinitely many. So a span is not one fixed shape -- it is a \emph{shape of a given dimension, oriented some particular way}. On your walk, try to feel the difference between ``what dimension is this span'' and ``which way is it tilted.'' Both questions matter; the first is about how many independent directions, the second about orientation.
\end{walkbox}

\section{A Fully Worked Span Question}

Definitions are cheap. Let us answer a real question by hand, because computing one span end-to-end teaches more than ten definitions.

\begin{example}[Is $\vec{b}$ in the span?]\label{ch03:ex-in-span}
Let
\[
  \vec{v}_1 = \colvec{1\\2}, \qquad \vec{v}_2 = \colvec{3\\1}, \qquad \vec{b} = \colvec{5\\0}.
\]
Is $\vec{b} \in \spn\{\vec{v}_1, \vec{v}_2\}$? In words: can we hit the target $\vec{b}$ by mixing $\vec{v}_1$ and $\vec{v}_2$?
\end{example}

Being in the span means there \emph{exist} weights $c_1, c_2$ with
\[
  c_1\colvec{1\\2} + c_2\colvec{3\\1} = \colvec{5\\0}.
\]
Reading off each coordinate turns this single vector equation into two scalar equations -- one per row:
\[
  \begin{aligned}
    c_1 + 3c_2 &= 5 \quad (\text{top row}),\\
    2c_1 + c_2 &= 0 \quad (\text{bottom row}).
  \end{aligned}
\]
From the second equation, $c_2 = -2c_1$. Substitute into the first:
\[
  c_1 + 3(-2c_1) = 5 \;\Longrightarrow\; c_1 - 6c_1 = 5 \;\Longrightarrow\; -5c_1 = 5 \;\Longrightarrow\; c_1 = -1.
\]
Then $c_2 = -2(-1) = 2$. Let us \emph{check} -- always check:
\[
  -1\colvec{1\\2} + 2\colvec{3\\1}
  = \colvec{-1\\-2} + \colvec{6\\2}
  = \colvec{5\\0} = \vec{b}. \checkmark
\]
Yes. $\vec{b}$ is in the span, reached with weights $(-1, 2)$. And because $\vec{v}_1, \vec{v}_2$ point in different directions, they span \emph{all} of $\R^2$ -- so this would have worked for any target $\vec{b}$ we picked.

\begin{focusbox}
Look closely at what just happened. ``Is $\vec{b}$ in the span?'' became ``does this system of linear equations have a solution?'' That is not a coincidence -- it is the bridge between geometry (spans, planes) and algebra (systems, elimination). Every time you solve a linear system you are secretly asking a span question. Keep this in your pocket; it is the heart of the next several chapters.
\end{focusbox}

\section{The Column Picture: $\mat{A}\vec{x}$ as a Combination of Columns}

We now reveal what is, in Strang's own words, the most important idea in the whole subject. It reframes matrix-times-vector so that it \emph{becomes} a linear combination. Once you see it you cannot unsee it.

Take a matrix and a vector:
\[
  \mat{A} = \begin{bmatrix} 1 & 3 \\ 2 & 1 \end{bmatrix}, \qquad \vec{x} = \colvec{c_1 \\ c_2}.
\]
The usual ``rows times the vector'' recipe for $\mat{A}\vec{x}$ gives
\[
  \mat{A}\vec{x}
  = \begin{bmatrix} 1 & 3 \\ 2 & 1 \end{bmatrix}\colvec{c_1\\c_2}
  = \colvec{1\cdot c_1 + 3\cdot c_2 \\ 2\cdot c_1 + 1\cdot c_2}.
\]
Now stare at that right-hand side and \emph{regroup it by columns} instead of by rows:
\[
  \colvec{c_1 + 3c_2 \\ 2c_1 + c_2}
  = \colvec{c_1 \\ 2c_1} + \colvec{3c_2 \\ c_2}
  = c_1\colvec{1\\2} + c_2\colvec{3\\1}.
\]
Look at what those two final vectors are: $\colvec{1\\2}$ is the first \emph{column} of $\mat{A}$, and $\colvec{3\\1}$ is the second column. So
\[
  \mat{A}\vec{x} = c_1\,(\text{column 1 of }\mat{A}) + c_2\,(\text{column 2 of }\mat{A}).
\]

\begin{formalbox}[The column rule for $\mat{A}\vec{x}$]
If $\mat{A}$ has columns $\vec{a}_1,\vec{a}_2,\dots,\vec{a}_n$ (each a vector in $\R^m$) and $\vec{x} = \colvec{x_1\\ \vdots \\ x_n}$, then
\[
  \mat{A}\vec{x} = x_1\vec{a}_1 + x_2\vec{a}_2 + \cdots + x_n\vec{a}_n.
\]
\textbf{$\mat{A}\vec{x}$ is the linear combination of the columns of $\mat{A}$ with weights taken from $\vec{x}$.}
\end{formalbox}

This is the column picture. The entries of $\vec{x}$ are not mysterious -- they are exactly the \emph{weights} in a linear combination, and the things being combined are the columns. Multiplying by a matrix \emph{is} taking a linear combination of its columns. Full stop.

Figure~\ref{ch03:column-picture} draws this for the system we solved by hand in Example~\ref{ch03:ex-in-span}: the columns of $\mat{A}$ are the two arrows you are allowed to scale, and $\mat{A}\vec{x}=\vec{b}$ just asks you to dial the weights until tip-to-tail you land exactly on the target $\vec{b}$.

\begin{figure}[H]
\centering
\begin{tikzpicture}[scale=0.7,>=Stealth]
  \draw[help lines,gray!22] (-2.4,-2.9) grid (6.4,2.9);
  \draw[->] (-2.4,0)--(6.4,0) node[right]{$x$};
  \draw[->] (0,-2.9)--(0,2.9) node[above]{$y$};
  % columns a1=(1,2), a2=(3,1)
  \draw[->,vecblue,very thick] (0,0)--(1,2) node[midway,left]{$\vec{a}_1$};
  \draw[->,vecred,very thick] (0,0)--(3,1) node[midway,below]{$\vec{a}_2$};
  % combination -1 a1 + 2 a2 = (5,0): walk -a1 then +2a2
  \draw[->,vecblue,thick,dashed] (0,0)--(-1,-2) node[midway,right]{$-\vec{a}_1$};
  \draw[->,vecred,thick,dashed] (-1,-2)--(5,0);
  \node[vecred!70!black] at (2.3,-1.5) {\small $2\vec{a}_2$};
  \draw[->,targetgreen,very thick] (0,0)--(5,0);
  \fill[targetgreen] (5,0) circle (2.8pt) node[above right,targetgreen!60!black]{\small $\vec{b}=\colvec{5\\0}$};
\end{tikzpicture}
\caption{The column picture of $\mat{A}\vec{x}=\vec{b}$ for $\mat{A}=\protect\begin{bmatrix}1&3\\2&1\protect\end{bmatrix}$, $\vec{b}=\colvec{5\\0}$. Solving the system means choosing weights so that $x_1\vec{a}_1+x_2\vec{a}_2$ reaches the green target. The weights $(-1,2)$ from Example~\ref{ch03:ex-in-span} do it: walk $-\vec{a}_1$ (down-left), then add $2\vec{a}_2$, and you land on $\vec{b}$.}
\label{ch03:column-picture}
\end{figure}

\begin{intuitionbox}
There are two ways to read $\mat{A}\vec{x}$, and a master keeps both in mind:
\begin{itemize}
  \item \textbf{Row picture:} each entry of the output is a dot product of a row with $\vec{x}$. Good for computing one number at a time, and for thinking of each equation as a line/plane.
  \item \textbf{Column picture:} the whole output is a linear combination of the columns, with weights $\vec{x}$. Good for thinking about \emph{which outputs are reachable at all}.
\end{itemize}
The row picture answers ``what is the answer?'' The column picture answers ``which answers are even possible?'' This chapter lives in the column picture.
\end{intuitionbox}

\begin{labbox}
Do this one with pen and paper before reading on. Let
\[
  \mat{A} = \begin{bmatrix} 2 & 0 & 1 \\ 1 & 3 & 1 \end{bmatrix}, \qquad \vec{x} = \colvec{1\\2\\-1}.
\]
Compute $\mat{A}\vec{x}$ \emph{twice}: once the row way (two dot products), once the column way ($1$ times column one, plus $2$ times column two, minus $1$ times column three, added up). You should get the same vector both times -- $\colvec{1\\6}$. Feeling the two routes land on the same point is the whole lesson.
\end{labbox}

\subsection{What the column picture says about solving $\mat{A}\vec{x} = \vec{b}$}

Here is the payoff. The equation $\mat{A}\vec{x} = \vec{b}$ asks: \emph{find weights $\vec{x}$ so that the linear combination of the columns equals $\vec{b}$}. In our column language:
\[
  \mat{A}\vec{x} = \vec{b}
  \quad\Longleftrightarrow\quad
  x_1\vec{a}_1 + \cdots + x_n\vec{a}_n = \vec{b}
  \quad\Longleftrightarrow\quad
  \vec{b} \in \spn\{\vec{a}_1,\dots,\vec{a}_n\}.
\]
Read that chain slowly. \textbf{The system $\mat{A}\vec{x}=\vec{b}$ has a solution if and only if $\vec{b}$ lies in the span of the columns of $\mat{A}$.} The span of the columns is so important it gets a name -- the \textbf{column space} of $\mat{A}$, written $\col(\mat{A})$ -- and an entire chapter later in the book. You are meeting it here as ``the set of all reachable right-hand sides.''

\begin{focusbox}
Solvability is a geometry question. ``Can I solve $\mat{A}\vec{x}=\vec{b}$?'' is \emph{exactly} ``is the target $\vec{b}$ sitting inside the span of the columns?'' If $\vec{b}$ is in the column space, a solution exists. If $\vec{b}$ floats outside it, no choice of weights will ever reach it, and the system is unsolvable -- not because you were not clever enough, but because the target is geometrically out of reach.
\end{focusbox}

\section{Reaching a Target $\vec{b}$}

Let us turn solvability into pictures with a few small cases, because the geometry tells you the answer before any algebra does.

\textbf{Case 1: columns span everything.} Take
\[
  \mat{A} = \begin{bmatrix} 1 & 3 \\ 2 & 1 \end{bmatrix}.
\]
Its columns $\colvec{1\\2}$ and $\colvec{3\\1}$ point in different directions, so they span all of $\R^2$. Therefore \emph{every} target $\vec{b}\in\R^2$ is reachable -- $\mat{A}\vec{x}=\vec{b}$ is solvable for every $\vec{b}$, exactly as we saw in Example~\ref{ch03:ex-in-span}. The column space fills the whole plane; nothing is out of reach.

\textbf{Case 2: columns collapse to a line.} Take
\[
  \mat{B} = \begin{bmatrix} 1 & 2 \\ 2 & 4 \end{bmatrix}.
\]
Its second column $\colvec{2\\4}$ is exactly twice the first column $\colvec{1\\2}$. They lie on the same line. So no matter what weights $x_1, x_2$ you choose,
\[
  x_1\colvec{1\\2} + x_2\colvec{2\\4} = (x_1 + 2x_2)\colvec{1\\2}
\]
is always a multiple of $\colvec{1\\2}$. The column space is a single line -- the line through the origin in direction $\colvec{1\\2}$. The system $\mat{B}\vec{x}=\vec{b}$ is solvable \emph{only} for targets $\vec{b}$ sitting on that line, and unsolvable for every $\vec{b}$ off it.

\begin{figure}[H]
\centering
\begin{tikzpicture}[scale=0.8,>=Stealth]
  \draw[help lines,gray!22] (-1.4,-1.4) grid (4.4,4.9);
  \draw[->] (-1.4,0)--(4.4,0) node[right]{$x$};
  \draw[->] (0,-1.4)--(0,4.9) node[above]{$y$};
  % column space line: direction (1,2)
  \draw[sourceblue!55,very thick] (-0.7,-1.4)--(2.45,4.9);
  \draw[->,vecblue,very thick] (0,0)--(1,2) node[left]{$\vec{a}_1$};
  \draw[->,vecred,very thick] (0,0)--(2,4) node[right]{$\vec{a}_2=2\vec{a}_1$};
  \node[sourceblue!70!black] at (2.9,4.2) {\small $\col(\mat{B})$};
  % reachable target
  \fill[targetgreen] (1.5,3) circle (2.6pt) node[right,targetgreen!60!black]{\small reachable $\vec{b}$};
  % unreachable target
  \fill[trapred] (3,1) circle (2.6pt) node[right,trapred]{\small unreachable $\vec{b}$};
\end{tikzpicture}
\caption{When the columns are parallel, the column space collapses to a line. Targets on the line (green) are reachable; targets off the line (red) are forever out of reach, no matter the weights.}
\label{ch03:collapse-line}
\end{figure}

\begin{example}[Reading solvability off the geometry]\label{ch03:ex-solvable}
For $\mat{B}=\begin{bmatrix}1&2\\2&4\end{bmatrix}$ above, is $\mat{B}\vec{x}=\colvec{3\\6}$ solvable? Is $\mat{B}\vec{x}=\colvec{3\\1}$ solvable?

\emph{First target.} $\colvec{3\\6} = 3\colvec{1\\2}$ sits right on the line, so yes -- pick $x_1=3, x_2=0$ (or $x_1=1,x_2=1$, or infinitely many others, since $x_1+2x_2=3$ is one equation in two unknowns). Reachable, in infinitely many ways.

\emph{Second target.} $\colvec{3\\1}$ is \emph{not} a multiple of $\colvec{1\\2}$ (that would need the second entry to be $6$, not $1$). It floats off the line. No weights reach it. Unsolvable. The algebra confirms it: the two equations $x_1+2x_2=3$ and $2x_1+4x_2=1$ say $2(x_1+2x_2)$ equals both $6$ and $1$ -- a flat contradiction.
\end{example}

\begin{trapbox}
A square matrix does \emph{not} guarantee a solution. Both $\mat{A}$ and $\mat{B}$ above are $2\times 2$, but $\mat{A}$ solves for every $\vec{b}$ while $\mat{B}$ fails for most. The difference is not the shape of the matrix -- it is whether the columns actually span the target space, or secretly collapse to something smaller. Always ask ``what do the columns span?'' before assuming a solution exists.
\end{trapbox}

\begin{applicationbox}
\textbf{Color, recipes, and recommender systems.} Linear combinations are how machines mix and reach things.

A screen pixel is a linear combination of three column vectors -- red, green, blue ``primaries.'' The weights $(r,g,b)$ are the dials. The span of those three columns is the set of \emph{all displayable colors} (the monitor's gamut). A color outside that span literally cannot be shown on the screen -- it is an unreachable $\vec{b}$, exactly like the red dot in Figure~\ref{ch03:collapse-line}. Wider-gamut displays add primaries that point in ``newer directions,'' enlarging the span.

The same idea powers recommender systems. In a model like matrix factorization, each user's taste is a short vector of weights, and a predicted rating is a linear combination of a handful of hidden ``taste direction'' vectors. The whole catalog of predictable preferences is the span of those directions. When Netflix or Spotify says ``because you liked X,'' it is reporting which combination of latent direction-vectors your taste landed on. The dimension of that span -- how many genuinely different taste directions the model keeps -- is a design knob (the ``rank''), and you will meet it precisely when we get to the SVD and PCA.
\end{applicationbox}

\section{When a New Vector Adds Nothing: Dependence, Previewed}

We keep bumping into the same phenomenon: sometimes a new vector enlarges the span, sometimes it is dead weight. Let us name the pattern, gently, because the full theory comes later but the geometry is fully visible now.

\begin{definition}[Adds nothing / redundant -- informal]\label{ch03:def-redundant}
A vector $\vec{v}$ is \textbf{redundant} for a collection if it already lies in the span of the others -- that is, $\vec{v}$ is itself a linear combination of them. Adding a redundant vector does not change the span.
\end{definition}

Two quick litmus tests you can run by eye:

\begin{itemize}
  \item \textbf{One vector is a multiple of another} ($\vec{v}_2 = c\,\vec{v}_1$): redundant. They share a line. This is the case of $\mat{B}$ above.
  \item \textbf{One vector lies in the plane of two others} ($\vec{v}_3 = a\vec{v}_1 + b\vec{v}_2$): redundant. It adds no new direction out of that plane.
\end{itemize}

When \emph{no} vector in the set is redundant -- when every one of them points somewhere the others cannot reach -- we call the set \textbf{linearly independent}. When at least one is redundant, the set is \textbf{linearly dependent}. That is the whole geometric content of those two intimidating words, and you understand it already.

\begin{theorem}[Two non-parallel vectors span all of $\R^2$]\label{ch03:thm-span-r2}
Let $\vec{v}_1,\vec{v}_2 \in \R^2$ be nonzero and non-parallel (neither is a scalar multiple of the other). Then $\spn\{\vec{v}_1,\vec{v}_2\} = \R^2$: every vector in the plane is some linear combination of them.
\end{theorem}

\begin{proof}
Write $\vec{v}_1 = \colvec{a\\c}$ and $\vec{v}_2 = \colvec{b\\d}$, and let $\vec{b}^\star = \colvec{p\\q}$ be \emph{any} target. We must find weights $x_1,x_2$ with $x_1\vec{v}_1 + x_2\vec{v}_2 = \vec{b}^\star$, i.e.
\[
  \begin{aligned}
    a x_1 + b x_2 &= p,\\
    c x_1 + d x_2 &= q.
  \end{aligned}
\]
Solve by elimination. Multiply the first equation by $d$, the second by $b$, and subtract:
\[
  (ad - bc)\,x_1 = dp - bq.
\]
The number $ad-bc$ is the \emph{determinant} of the matrix with columns $\vec{v}_1,\vec{v}_2$. The non-parallel hypothesis is exactly the statement $ad-bc \neq 0$ (parallel columns make it zero -- a fact we will prove in the determinant chapter, and which you can check on $\mat{B}$: $1\cdot 4 - 2\cdot 2 = 0$). Since $ad-bc\neq 0$ we may divide:
\[
  x_1 = \frac{dp - bq}{ad-bc}, \qquad\text{and similarly}\qquad x_2 = \frac{aq - cp}{ad-bc}.
\]
These weights exist for \emph{every} $(p,q)$, so every target is reachable. Hence the span is all of $\R^2$.
\end{proof}

\begin{remark}
The proof quietly did something lovely: it showed that ``non-parallel'' (a geometric statement) is the same as ``$ad-bc\neq 0$'' (an algebraic one), which is the same as ``the columns span $\R^2$'' (a span statement), which is the same as ``$\mat{A}\vec{x}=\vec{b}$ always solvable'' (a system statement). Four sentences, one truth, four dialects. Learning to translate between these dialects fluently \emph{is} learning linear algebra.
\end{remark}

\subsection{Counting: why you cannot span $\R^3$ with two vectors}

Here is a fact your geometric intuition already insists on, and it is worth saying out loud.

\begin{proposition}[Too few vectors cannot span]\label{ch03:prop-too-few}
The span of $k$ vectors in $\R^n$ is at most $k$-dimensional. In particular, fewer than $n$ vectors can never span all of $\R^n$.
\end{proposition}

We will only sketch the reasoning now (the airtight proof waits for the dimension chapter). Each vector you add can push the span up by \emph{at most} one dimension -- one new direction at most. Starting from nothing (dimension $0$), after $k$ vectors you have climbed at most $k$ rungs of the point-line-plane-space ladder. So two vectors max out at a plane; they can decorate $\R^3$ but never fill it. You need at least three vectors to span $\R^3$, at least $n$ to span $\R^n$. Conversely, more vectors than necessary guarantees redundancy: any $n+1$ vectors in $\R^n$ \emph{must} be dependent, because the $(n{+}1)$-th has nowhere new to point.

\begin{connectionbox}
\textbf{Spans are everywhere once you know the shape.} The pattern ``set of all combinations of a few generators'' reaches far past arrows in $\R^n$.
\begin{itemize}
  \item \textbf{Polynomials.} The functions $1, t, t^2$ are ``vectors,'' and their span is every quadratic $a + bt + ct^2$. Calculus students compute in this span constantly without naming it. Taylor series live in the span of $1, t, t^2, t^3, \dots$
  \item \textbf{Differential equations.} The solutions of $y'' + y = 0$ are exactly the linear combinations $c_1\cos t + c_2\sin t$ -- a $2$-dimensional span of two ``basis solutions.'' Solving a linear ODE is finding generators and taking their span.
  \item \textbf{Signals and Fourier.} A sound wave is (approximately) a linear combination of pure sine tones. The span of the basic tones is the space of signals you can build; an equalizer is just you, turning the weight dials.
  \item \textbf{Geometry of subspaces.} ``Line through the origin,'' ``plane through the origin,'' and ``all of space'' are not three unrelated shapes -- they are spans of $1$, $2$, and $3$ independent vectors. Span is the engine that manufactures \emph{subspaces}, the topic that organizes the middle of this book.
\end{itemize}
Whenever you see ``all combinations of a few building blocks,'' you are looking at a span, and everything in this chapter applies.
\end{connectionbox}

\begin{strangbox}
\textbf{Before Strang:} when he opens his course with the ``column picture'' and insists $\mat{A}\vec{x}$ is a combination of columns, you will already own it -- you derived it above by regrouping $\colvec{c_1+3c_2\\2c_1+c_2}$ into columns. His famous line, ``linear algebra is about $\mat{A}\vec{x}=\vec{b}$, and the question is whether $\vec{b}$ is a combination of the columns,'' is precisely the chain we built in the section on reaching a target.

\textbf{After Strang:} when you return for the \emph{column space} $\col(\mat{A})$ and the \emph{four fundamental subspaces}, recognize that the column space is nothing more than the span you met here, now given its formal home. The rank of a matrix is the dimension of this span -- how many rungs of the ladder its columns actually climb. Every ``rank'' theorem later is a ``how big is the span'' theorem in disguise.
\end{strangbox}

\begin{aestheticbox}
Step back and admire the economy. Two humble moves -- scale and add -- generate the entire geometric vocabulary of the subject: lines, planes, spaces, and the precise sense in which one vector can be ``new'' or ``already implied.'' From those moves alone we got a clean test for when an equation is solvable, a dictionary connecting geometry to algebra to systems, and a preview of independence, rank, and column space. Almost nothing went in; an enormous amount came out. That leverage -- tiny rules, vast structure -- is what makes linear algebra feel less like a list of tricks and more like a landscape you can learn to see.
\end{aestheticbox}

\section{Summary}

\begin{itemize}
  \item A \textbf{linear combination} $c_1\vec{v}_1 + \cdots + c_k\vec{v}_k$ is built from only two moves: \emph{scaling} (along a direction) and \emph{adding} (mixing directions). The $c_i$ are weights.
  \item The \textbf{span} is the set of \emph{all} linear combinations -- every destination the weights can reach. It always passes through the origin.
  \item Geometrically the span climbs a ladder: \emph{point} (zero vectors), \emph{line} (one direction), \emph{plane} (two independent directions), \emph{space}, and on up. Each new vector adds at most one rung.
  \item A vector that already lies in the others' span is \textbf{redundant}; it adds nothing. ``No redundancy'' is \textbf{linear independence}; this is the geometric seed of independence, rank, and dimension.
  \item \textbf{Column picture:} $\mat{A}\vec{x} = x_1\vec{a}_1 + \cdots + x_n\vec{a}_n$ is a linear combination of the columns of $\mat{A}$ with weights $\vec{x}$.
  \item Therefore $\mat{A}\vec{x}=\vec{b}$ is solvable \emph{exactly when} $\vec{b}$ lies in the span of the columns (the \textbf{column space}). Solvability is a geometry question.
\end{itemize}

\begin{recoverybox}
\textbf{Cost and intuition recap.} The \emph{intuition} is a single sentence: a span is everywhere you can get to by scaling and adding, and it always grows by clean rungs -- point, line, plane, space. The \emph{cost} is just as small: to test whether a target $\vec{b}$ lies in $\spn\{\vec{v}_1,\dots,\vec{v}_k\}$, you do not search forever -- you set up one linear system, $c_1\vec{v}_1+\cdots+c_k\vec{v}_k=\vec{b}$ ($n$ equations, $k$ unknowns), and solve it once. ``In the span?'' costs exactly ``does this system have a solution?'' -- no more. Hold those two together: the picture (rungs of a ladder) tells you \emph{what} a span is; the system (solve once, then check) tells you \emph{how much} it costs to use one.
\end{recoverybox}

\section{Exercises}

\begin{exercise}[Warmup combinations]\label{ch03:ex-warmup}
Let $\vec{u} = \colvec{1\\-1}$ and $\vec{w} = \colvec{2\\3}$. Compute each linear combination:
(a) $3\vec{u} + \vec{w}$;\;
(b) $2\vec{u} - 2\vec{w}$;\;
(c) $0\vec{u} + 5\vec{w}$;\;
(d) $-\vec{u} - \vec{w}$.
For (c), what does the weight $0$ on $\vec{u}$ tell you about where the result lands?
\end{exercise}

\begin{exercise}[Solve for the weights]\label{ch03:ex-find-weights}
Find weights $c_1, c_2$ so that $c_1\colvec{1\\1} + c_2\colvec{1\\-1} = \colvec{4\\2}$. Then verify your answer by substituting back. (Hint: add and subtract the two scalar equations.)
\end{exercise}

\begin{exercise}[In the span or not?]\label{ch03:ex-membership}
Let $\vec{v}_1 = \colvec{2\\1}$ and $\vec{v}_2 = \colvec{-4\\-2}$.
(a) What is $\spn\{\vec{v}_1,\vec{v}_2\}$ geometrically -- a point, a line, or a plane? Why?
(b) Is $\colvec{6\\3}$ in the span? Is $\colvec{1\\0}$ in the span? Justify each answer geometrically, then confirm algebraically.
\end{exercise}

\begin{exercise}[Column picture by hand]\label{ch03:ex-column}
Let $\mat{A} = \begin{bmatrix} 1 & 0 & 2 \\ 0 & 1 & 3 \end{bmatrix}$ and $\vec{x} = \colvec{4\\-1\\2}$. Compute $\mat{A}\vec{x}$ as a linear combination of the columns of $\mat{A}$ (write out ``$4$ times column one, plus \dots''). Then state what $\spn$ of the three columns is, and explain why the third column is redundant.
\end{exercise}

\begin{exercise}[Reaching a target]\label{ch03:ex-target}
For $\mat{A} = \begin{bmatrix} 1 & 2 \\ 3 & 6 \end{bmatrix}$:
(a) Describe $\col(\mat{A})$ (the span of the columns) as a specific line through the origin.
(b) Is $\mat{A}\vec{x} = \colvec{2\\6}$ solvable? Is $\mat{A}\vec{x} = \colvec{2\\5}$ solvable? Answer by checking whether each target lies on the line from (a).
(c) For the solvable case, find \emph{all} solutions $\vec{x}$, and explain geometrically why there are infinitely many.
\end{exercise}

\begin{exercise}[The third vector decides]\label{ch03:ex-third}
In $\R^3$, let $\vec{v}_1 = \colvec{1\\0\\0}$ and $\vec{v}_2 = \colvec{0\\1\\0}$.
(a) Describe $\spn\{\vec{v}_1,\vec{v}_2\}$ as a specific plane.
(b) For $\vec{v}_3 = \colvec{2\\3\\0}$, what is $\spn\{\vec{v}_1,\vec{v}_2,\vec{v}_3\}$? Did $\vec{v}_3$ add a new direction?
(c) For $\vec{v}_3 = \colvec{2\\3\\1}$, what is $\spn\{\vec{v}_1,\vec{v}_2,\vec{v}_3\}$ now? What changed?
\end{exercise}

\begin{exercise}[A span that misses the origin?]\label{ch03:ex-affine}
Your friend claims the line $y = x + 1$ in the plane is the span of some single vector. Prove your friend is wrong in one sentence. (Hint: which point is in \emph{every} span, and is it on this line?)
\end{exercise}

\begin{exercise}[Counting argument]\label{ch03:ex-counting}
Explain, in your own words and with a picture, why three vectors in $\R^2$ \emph{must} be linearly dependent (at least one is redundant). Then exhibit explicit weights, not all zero, with $c_1\colvec{1\\0} + c_2\colvec{0\\1} + c_3\colvec{1\\1} = \colvec{0\\0}$, demonstrating the dependence concretely.
\end{exercise}

\begin{exercise}[Same span, different generators]\label{ch03:ex-same-span}
Show that $\spn\left\{\colvec{1\\0},\colvec{0\\1}\right\} = \spn\left\{\colvec{1\\1},\colvec{1\\-1}\right\}$ by proving each side contains the other. (To show ``$\subseteq$,'' write a general vector of one span as a combination of the other set's vectors.) What does this say about the idea that a span has many possible ``generating sets''?
\end{exercise}

\begin{exercise}[Proof: span is closed under combination]\label{ch03:ex-closed}
Let $S = \spn\{\vec{v}_1,\dots,\vec{v}_k\}$. Suppose $\vec{p}$ and $\vec{q}$ are both in $S$. Prove that $\alpha\vec{p} + \beta\vec{q}$ is also in $S$ for any scalars $\alpha,\beta$. (Write $\vec{p}$ and $\vec{q}$ as combinations of the $\vec{v}_i$, then combine.) This ``closed under linear combination'' property is the defining feature of a \emph{subspace} -- you have just proved every span is one.
\end{exercise}

\begin{exercise}[Challenge: build any vector two ways]\label{ch03:ex-challenge}
Let $\vec{v}_1 = \colvec{1\\0}$, $\vec{v}_2 = \colvec{0\\1}$, $\vec{v}_3 = \colvec{1\\1}$.
(a) Find two \emph{different} sets of weights $(c_1,c_2,c_3)$ that both produce $\colvec{3\\2}$.
(b) Explain why having a redundant vector in the set is exactly what makes more than one set of weights possible. Connect this to why $\mat{B}\vec{x}=\colvec{3\\6}$ in Example~\ref{ch03:ex-solvable} had infinitely many solutions.
\end{exercise}
